\documentclass{article}
\usepackage[T1]{fontenc}
\usepackage{textcomp}
\usepackage{newtxtext,newtxmath}

\usepackage{hyperref}
\usepackage{tabularx}
\usepackage{xurl}
\usepackage{graphicx}


\usepackage {tikz}
\usepackage[utf8]{inputenc}
\usepackage[mono=false]{libertinus}
\usepackage[T1]{fontenc}
\usepackage{sectsty}
\allsectionsfont{\sffamily}
\usepackage{tocloft}
\usepackage{multicol}
\renewcommand{\cftsecfont}{\bfseries\sffamily}
\usepackage[numbers]{natbib}
\usepackage{enumitem,lipsum}
\usepackage{graphicx}
\usepackage{float}
\usepackage{subcaption}
\usepackage{geometry}
\geometry{margin=0.9in}
\usepackage{amsmath}
\usepackage{amssymb}
\usepackage{amsthm}
\usepackage{fancyhdr}
\usepackage{array}
\usepackage[table,xcdraw]{xcolor}
\usepackage{hyperref}
\usepackage[breakable]{tcolorbox}
\usepackage{bbm}
\usepackage{parskip}
\usepackage{physics}
\usepackage[hang,flushmargin]{footmisc}
\usepackage{mathtools}
\usepackage{stmaryrd}
\usepackage{tikz}
\usepackage{pgfplots}
\pgfplotsset{compat=1.15}
% \usepackage{mathrsfs}
\usetikzlibrary{arrows}
\usepackage{comment}
\usepackage{blkarray}
\usepackage{multirow}
\begin{document}

\section{Motivation}
Precipitation remains one of the least predictable components of the Earth system. This difficulty shows up across modeling paradigms: traditional climate and NWP models struggle with drizzle bias, phase errors, and extremes; modern AI weather systems often inherit similar issues despite high skill on large-scale flow. The reasons are well known: rain is intermittent and thresholded, triggered by strongly nonlinear interactions across scales, and sensitive to small errors in humidity, stability, and vertical motion. Microphysical and convective processes are only partially resolved on typical grids, creating intrinsic stochasticity at the resolved scale.

Most AI forecasts advance the atmospheric \emph{state} autoregressively in time and then obtain precipitation either (i) by diagnosing it (Pangu)with a separate network from predicted state variables, or (ii) by predicting precipitation jointly with other targets (GenCast), typically at 3–6h time steps. This design leaves a basic scientific question under-explored: \emph{even if the state variables were predicted perfectly, how far into the future is precipitation itself predictable?}

We propose to quantify an \textbf{intrinsic window of predictability for precipitation} that is \emph{conditional on perfect knowledge of the large-scale state}. Formally, let $X_t$ denote a set of resolved prognostic variables (e.g., ${T,q,\mathbf{u},\omega,  \ldots}$) and let $P_t$ denote precipitation. We study the conditional distribution
\[p(P(t+\tau)| X(t:t+m))\]

and define a lead-time–dependent skill function $\mathcal{S}(\tau)$ (e.g., CRPS skill, Brier skill, or mutual information $I(P_{t+\tau};X_{t:t+m})$). The \emph{predictability horizon} $\tau^\star$ is the smallest $\tau$ for which $\mathcal{S}(\tau)$ drops below a pre-specified threshold (near persistence/climatology). This horizon represents an upper bound: no autoregressive AI system that relies on these state variables can exceed it, because it reflects uncertainty in mapping state to rain rather than errors in state prediction.

By separating \emph{state-evolution skill} from \emph{state-to-precipitation skill}, we (i) isolate where AI forecasts hit a hard ceiling, (ii) diagnose which variables (and temporal context $m$--- we can consider this in the future!) carry the most predictive information for rain occurrence, rate, and extremes, and (iii) assess how sampling cadence (0.25h vs.\ 3h) scale modify $\tau^\star$. The outcome is a principled target for future AI architectures: if the observed ceiling is low for a given regime or scale, progress requires adding new predictors (e.g., higher-resolution vertical structure or convective proxies) or embracing stochastic outputs, rather than further tuning the autoregressive core.



\section{Lead-Conditioned U-Net for Rain Occurrence and Amount}

We train a convolutional encoder--decoder (U-Net) that ingests large-scale state at time \(t\)---sea surface temperature (SST), column-integrated water vapor (PW), and outgoing longwave radiation (OLR)---and predicts precipitation at a fixed lead \(t+\tau\). Because precipitation fields are sparse and thresholded, we decompose the task into (i) \emph{occurrence} (rain/no-rain) and (ii) \emph{amount} (conditional intensity when it rains), sharing a single backbone and using two task heads. Formally, with \(P_{t+\tau}\) the true rate and threshold \(r_0\) (e.g., \(0.01\,\mathrm{mm\,h^{-1}}\)), define the binary label
\[
Z_{t+\tau} \;=\; \mathbb{1}\!\left[P_{t+\tau}\ge r_0\right],
\qquad
Y_{t+\tau} \;=\; P_{t+\tau}\, .
\]
For each lead \(\tau\), a separate model is trained to map \((\mathrm{X(t)})\mapsto (\widehat{p}_{t+\tau}, \widehat{Y}_{t+\tau})\), where \(\widehat{p}\in(0,1)\) is the occurrence probability and \(\widehat{Y}\ge 0\) is the amount.

\paragraph{Data.}
The dataset analyzed in this study is the DYAMOND MPAS-Atmosphere (MPAS) global cloud-resolving simulation, performed at ~3.8 km horizontal resolution with 15-min output over a 40-day winter period. MPAS employs a fully compressible, nonhydrostatic dynamical core on an unstructured mesh, enabling explicit simulation of convective processes. Cloud microphysics are represented with a double-moment scheme for cloud and rain hydrometeors. Global cloud-resolving models (GCRMs) explicitly resolve deep convection and cloud systems, providing dynamically consistent four-dimensional fields essential for process understanding and data-driven modeling (Satoh et al., 2019). However, GCRMs remain model-dependent and exhibit systematic biases. MPAS generally reproduces mean precipitation and its diurnal cycle within the multi-model range (Stevens et al., 2019) but tends to overestimate deep convective clouds and heavy rainfall compared with the IMERG precipitation product (0.1° × 0.1°, 30-min; Feng et al., 2025). Despite these limitations, the DYAMOND MPAS data provide high temporal resolution, physically consistent fields, and a uniform data structure, making them well suited for precipitation analysis and machine-learning applications.


\paragraph{Architecture.}
The network is a standard U-Net with skip connections and bilinear upsampling. ADD INFORMATION ABOUT THE NETWORK: an occurrence head outputs logits whose sigmoid gives \(\widehat{p}\); an amount head outputs parameters of a positive distribution for wet pixels. Add what the regressor does: is it predicting P directly or log P!

\paragraph{Training objective.}
TODO: what is our objective function?

\paragraph{Metrics and skill as a function of lead.}
TODO: We need to define our metrics!

Let \(\mathcal{S}_{\text{occ}}(\tau)\) and \(\mathcal{S}_{\text{amt}}(\tau)\) denote chosen skill measures (e.g., PR-AUC for occurrence and \(1-\)normalized CRPS for amounts) relative to a baseline. We summarize the decay with lead by fitting an exponential envelope \( \mathcal{S}(\tau)\approx \mathcal{S}_0 e^{-\tau/T_e}\), yielding e-folding times \(T_e^{\text{occ}}\) and \(T_e^{\text{amt}}\), and define a predictability horizon
\[
\tau^\star \;=\; \inf\{\tau:\ \mathcal{S}(\tau)\le \mathcal{S}_{\text{base}}+\epsilon\},
\]
where \(\mathcal{S}_{\text{base}}\) is persistence or climatology skill and \(\epsilon\) a small tolerance. Baselines include (i) spatially varying climatology (seasonal rain frequency and conditional mean) and (ii) persistence \(P_{t+\tau}=P_t\) (with occurrence from \(Z_t\) and amount from \(Y_t\) at wet pixels, zero otherwise), ensuring that \(\mathcal{S}(\tau)\) quantifies genuine lead information beyond trivial heuristics.

\paragraph{Scientific readouts.}
This design directly exposes how lead \(\tau\) degrades (a) the ability to identify where it will rain and (b) the fidelity of amounts where it does, conditional on perfect large-scale state at time \(t\). Comparing \(\mathcal{S}_{\text{occ}}(\tau)\) and \(\mathcal{S}_{\text{amt}}(\tau)\) reveals whether localization or intensity saturates first; stratifying the same curves by PW/OLR/SST anomalies (similar to SST-PW plot for example) links the observed horizons to physically interpretable regimes.


\section{Input Configurations \texorpdfstring{$\mathbf{X}(t)$}{X(t)}}

We evaluate how different choices of input state \(X(t)\) affect precipitation predictability at lead \(t+\tau\). The experiments are grouped into three progressively enriched configurations:

\paragraph{Case 1: Large-scale radiative--moist background.}
Inputs include precipitable water (PW), sea-surface temperature (SST), and outgoing longwave radiation (OLR) at time \(t\).
This baseline probes how far a coarse thermodynamic and radiative state---without explicit vertical structure---can carry precipitation occurrence and amount skill.

\paragraph{Case 2: Role of thermodynamics (vertical structure).}
We augment Case~1 with specific humidity \(q\) and temperature \(T\) at 850, 600, and 400\,hPa.
This configuration tests whether adding column stratification (moisture and stability) extends the predictability horizon by sharpening information about saturation deficit, lapse rate, and convective inhibition/available potential energy proxies.

\paragraph{Case 3: Role of dynamics.}
We further augment Case~2 with vertical velocity \(w\) (or \(\omega\)) at 850 and 400\,hPa, providing explicit large-scale ascent/descent as a trigger for convection and stratiform organization.

\paragraph{Coupling and interpretability.}
Thermodynamics and dynamics are strongly coupled: persistent moisture and stability anomalies co-vary with ascent patterns, so a substantial fraction of the information in \(w\) may already be encoded in \((q,T)\).
Our analysis explicitly asks whether providing \(w\) \emph{adds incremental predictive information} beyond thermodynamic structure: if skill-vs-lead curves \(\mathcal{S}(\tau)\) and their horizons \(\tau^\star\) improve when \(w\) is included, we gain evidence that explicit dynamics convey flow information not recoverable from thermodynamic proxies alone.


\section{Some thoughts based on what we have learned so far:}


There are two separable things our experiments are revealing:

1. What determines skill at t=0? — input variable sensitivity, answered by your variable ablation \\
2. What determines $\tau^\star$?  the intrinsic decorrelation timescale of precipitation-relevant atmospheric modes, independent of input choice


One interesting thing to look at is how $\tau^\star$ varies with location. It would be interesting to see if there is a difference between land vs ocean, and between tropics vs extra-tropics. 


Precipitation prediction can be conceptually decomposed into two distinct challenges. The first is predicting the \textit{amplitude} of convective events --- their intensity and vertical structure. The second, and more elusive, is predicting the \textit{trigger}: whether and where convection actually initiates. The large-scale thermodynamic environment, captured by variables such as column water vapor, outgoing longwave radiation, and moist static energy profiles, largely controls the amplitude. Given that convection occurs, how intense will it be? This component is relatively predictable because it inherits the memory timescale of the large-scale circulation. The trigger, by contrast, is governed by subgrid processes --- boundary layer thermals, cold pool collisions, and mesoscale convergence features --- that are effectively stochastic at the resolved scale of $0.25^\circ$. Even with perfect knowledge of the thermodynamic environment, the precise location and timing of convective initiation cannot be determined deterministically, introducing an irreducible uncertainty that grows with lead time.

It is important to note, however, that at zero lead time this trigger uncertainty is effectively eliminated by the input variables themselves. OLR is a near-direct observation of current convective state: deep convection emits at low brightness temperatures, clear sky at high. Similarly, PW tells the model where the atmospheric column is close to saturation. At $t = 0$, the model is not predicting the trigger --- it is reading it directly from the inputs. The residual error at zero lead time therefore reflects amplitude uncertainty, namely subgrid variability in convective intensity that cannot be resolved at $0.25^\circ$, rather than trigger stochasticity. As lead time increases, OLR and PW at $t_0$ progressively lose their ability to constrain where convection will fire at $t_0 + t$, and trigger uncertainty switches on. The skill decay we observe is precisely the timescale over which this transfer happens.

This decomposition can be made quantitative, and doing so carefully will significantly strengthen the interpretation of the results. The total prediction error at lead time $t$ can be written schematically as
%
\begin{equation}
    \epsilon(t) = \epsilon_A(t) + \epsilon_{\mathcal{T}}(t),
\end{equation}
%
[note before continue: the equation above assumes an additive decomposition. We can start by this and then modify it based on what we observe. the contribution might be multiplicative! ]
where $\epsilon_A$ captures amplitude uncertainty and $\epsilon_{\mathcal{T}}$ captures trigger uncertainty. To separate these empirically, we can proceed in two steps. First, we binarize precipitation at a physically meaningful threshold, say $P > 0.1$ mm~hr$^{-1}$ to distinguish raining from non-raining grid points (I believe you have done this already), and compute a categorical skill score for occurrence prediction as a function of lead time. Good choices are the Equitable Threat Score (ETS) or the Brier Skill Score (BSS), both of which measure how well the model identifies where convection fires relative to a climatological baseline, without any regard for intensity. The model should have a very good skill at time lag zero since we re shoing it the OLR and it tells everything about where convection s happening.The decay of this occurrence skill with lead time gives us an empirical handle on $\epsilon_{\mathcal{T}}(t)$. Second, we restrict verification to only those grid points where precipitation actually occurs in the target field, conditioning on observed occurrence, and compute $R^2$ or RMSE for intensity alone at those points. The decay of this conditioned skill is our empirical $\epsilon_A(t)$. If, as theory predicts, $\epsilon_{\mathcal{T}}$ grows faster than $\epsilon_A$, we have clean empirical evidence that trigger unpredictability dominates the predictability horizon. It would be also very interesting to see how they drop: linearly or exponentioally ? 

To make this even more precise, we fit an exponential model
%
\begin{equation}
    S(t) = \left(S_0 - S_\infty\right) e^{-t/\tau} + S_\infty
\end{equation}
%
to each skill curve separately, where $S_0$ is the skill at zero lead time, $S_\infty$ is the climatological floor, and $\tau$ is the e-folding decay timescale. This gives three interpretable parameters per experiment and per decomposition component. The story across the input variable experiments then becomes sharp and clean: adding thermodynamic profiles increases $S_0$ modestly, adding winds increases $S_0$ more substantially (we need to do the test with three baseline input + dynamics), but does $\tau$ remains essentially constant across all experiments?. The starting hypothesis is that the rate at which the atmosphere forgets is set by the intrinsic decorrelation timescale of the trigger, and no choice of input variables at $t_0$ can change that.
\end{document}